\documentclass[10pt,aspectratio=169]{beamer}
\usepackage[utf8]{inputenc}
\usepackage{listings}
\usepackage{hyperref}
\usepackage{xcolor}
\usepackage{ragged2e}
\usepackage{graphicx}
\usepackage{tikz}
\usetikzlibrary{positioning,calc,decorations.pathreplacing}

\usetheme{Dresden}
%\usetheme{Bergen}

\title{Programming Techniques 2026--2027}
\author{Hannu Parviainen}
\institute{Universidad de la Laguna}
\date{\today}

% Single spaces after full stops, so that e.g. \texttt{a .dot. b} has even spacing
\frenchspacing

\lstset{
  language=fortran,
  basicstyle=\ttfamily\scriptsize,
  columns=fullflexible,
  keepspaces=true,
  backgroundcolor=\color{white},
  keywordstyle=\color{blue},
  commentstyle=\color{gray},
  stringstyle=\color{red},
  showstringspaces=false,
  breaklines=true,
  frame=lines
}

\subtitle{Lecture 6: Derived Types and Custom Operators}

\begin{document}

\begin{frame}
  \titlepage
\end{frame}

\section{Derived Types}

\begin{frame}
	\centering
	\Huge \textbf{Derived Types}
\end{frame}

\begin{frame}[fragile]{Derived Types}

\begin{columns}[T]
  \column{0.5\textwidth}
  \begin{block}{Example: 3D vector type}
    \lstinputlisting[firstline=1,lastline=8]{ex_types.f90}
  \end{block}
  \begin{itemize}
  \item Basic Fortran types can be combined to create more complex \textit{derived} types.
  \item Derived type definitions should be placed in a module.
  \end{itemize}

  \column{0.5\textwidth}
  \begin{block}{Example: type usage}
  \lstinputlisting[firstline=10,lastline=16]{ex_types.f90}
  \end{block}
  \begin{itemize}
  \item Variables are declared with \texttt{type(name)}.
  \item A value is built with a constructor, which takes the components in the order of their definition.
  \end{itemize}

\end{columns}
\end{frame}

\begin{frame}[fragile]{Nested derived types}
  Previously defined types can be used as components of other derived types.
\begin{columns}[T,onlytextwidth]
  \column{0.5\textwidth}
  \begin{block}{Example: sphere}
    \lstinputlisting[firstline=1,lastline=12]{ex_types_2.f90}
  \end{block}

    \column{0.49\textwidth}
    \begin{block}{Example: sphere usage}
    \lstinputlisting[firstline=14,lastline=20,basicstyle=\ttfamily\scriptsize]{ex_types_2.f90}
  \end{block}
  \end{columns}
\end{frame}


% --- Derived Type Assignment ---
\begin{frame}[fragile]{Derived type assignment}
  \begin{columns}[T]
    \column{0.52\textwidth}
    \begin{block}{Ways to assign}
      \begin{itemize}
        \item Component-by-component using \texttt{\%}.
        \item As a whole object using a constructor.
        \item Assignment between two objects of the same derived type is intrinsic.
      \end{itemize}
      Inside a constructor, a derived-type component such as \texttt{centre} is given as a value of its own type: a variable or another constructor.
    \end{block}

    \column{0.48\textwidth}
\begin{lstlisting}[language=Fortran]
type(vector3d) :: pt1
type(sphere) :: p, ball

! component by component
pt1%x = 1.0
p%radius = 3.0
p%centre%x = 1.0

! whole object with constructors
pt1 = vector3d(1.0, 2.0, 3.0)
p%centre = vector3d(1.0, 2.0, 3.0)
p = sphere(p%centre, 10.0)
p = sphere(vector3d(1.0, 2.0, 3.0), 10.0)

! intrinsic assignment between objects
ball = p
\end{lstlisting}
  \end{columns}
\end{frame}

\begin{frame}[fragile]{Derived type I/O}
  \begin{columns}[T]
    \column{0.52\textwidth}
    \begin{block}{List-directed I/O}
      \begin{itemize}
        \item A derived type can be printed and read with \texttt{print *} and \texttt{read *} if it has no pointer or allocatable components.
        \item I/O proceeds component-by-component in order.
      \end{itemize}
      The two statements on the right print the same values.
    \end{block}

    \column{0.48\textwidth}
\begin{lstlisting}[language=Fortran]
print *, p

print *, p%centre%x, p%centre%y, &
         p%centre%z, p%radius
\end{lstlisting}
  \end{columns}
\end{frame}

\begin{frame}[fragile]{Allocatable and pointer components}
  \begin{columns}[T]
    \begin{column}{0.45\textwidth}
      \begin{itemize}
        \item A component can be allocatable (since Fortran 2003) or a pointer.
        \item Prefer allocatable components. They are freed automatically, and assignment copies their data (a \emph{deep copy}).
        \item A pointer component is only pointer-assigned in a copy, so both objects then share one target.
        \item Use pointer components for linked structures, such as a list node that points to the next node.
      \end{itemize}
    \end{column}

    \begin{column}{0.55\textwidth}
      \begin{lstlisting}
program components
  implicit none
  type :: spectrum
    character(len=:), allocatable :: name
    real, allocatable :: flux(:)
  end type spectrum
  type(spectrum) :: a, b
  a = spectrum('Vega', [1.0, 2.0, 3.0])
  b = a                  ! deep copy
  b%flux(1) = 99.0
  print *, a%flux        ! 1.0  2.0  3.0
  print *, b%flux        ! 99.0  2.0  3.0
end program components
      \end{lstlisting}
    \end{column}
  \end{columns}
\end{frame}


% --- Linked lists: drawing helpers ---
% A cell is drawn as a value box and a pointer box with a dot in it.
\tikzset{
  llval/.style={draw, fill=blue!15, minimum width=0.7cm, minimum height=0.55cm, inner sep=0pt, font=\ttfamily\small},
  llnext/.style={draw, fill=orange!25, minimum width=0.45cm, minimum height=0.55cm, inner sep=0pt},
  llptr/.style={->, >=stealth, thick},
  llvar/.style={font=\ttfamily\scriptsize, inner sep=1.5pt},
}
% \llcell[style]{name}{position}{value} creates the nodes (name)v and (name)n
\newcommand{\llcell}[4][]{%
  \node[llval, #1] (#2v) at (#3) {#4};
  \node[llnext, #1, anchor=west] (#2n) at (#2v.east) {};
  \fill (#2n.center) circle (1.3pt);
}

\begin{frame}[fragile]{Linked lists}
  \centering
  \begin{tikzpicture}
    \node[llnext, label={[llvar]above:head}] (h) at (0,0) {};
    \fill (h.center) circle (1.3pt);
    \llcell{a}{1.7,0}{1}
    \llcell{b}{3.8,0}{2}
    \llcell{c}{5.9,0}{3}
    \node[llvar, anchor=west] (nil) at (7.6,0) {null()};
    \draw[llptr] (h.center) -- (av.west);
    \draw[llptr] (an.center) -- (bv.west);
    \draw[llptr] (bn.center) -- (cv.west);
    \draw[llptr] (cn.center) -- (nil.west);
    \node[llvar, below] at (av.south) {val};
    \node[llvar, below] at (an.south) {next};
  \end{tikzpicture}

  \vspace{0.5em}
  \begin{columns}[T]
    \begin{column}{0.5\textwidth}
      \begin{itemize}
        \item A linked list is a chain of cells. Each cell holds a value and a pointer to the next cell.
        \item \texttt{head} points to the first cell. The pointer of the last cell is null.
        \item Every cell is allocated separately, so the list grows and shrinks one cell at a time, without copying.
        \item The price: there is no direct access to cell number $i$. You have to follow the pointers.
      \end{itemize}
    \end{column}

    \begin{column}{0.5\textwidth}
      \lstinputlisting[firstline=3,lastline=6]{linked_list.f90}
      Adding cells at the front builds the list above:
      \lstinputlisting[firstline=17,lastline=22]{linked_list.f90}
    \end{column}
  \end{columns}
\end{frame}


\begin{frame}[fragile]{Linked lists: inserting a cell}
  \begin{columns}[T]
    \begin{column}{0.48\textwidth}
      \centering
      \begin{tikzpicture}
        % Step 1: the new cell points at the rest of the list
        \begin{scope}
          \node[font=\scriptsize\bfseries, anchor=west] at (-0.5,1.25) {1.\ \texttt{new\%next => cur\%next}};
          \llcell{a}{0,0}{1}
          \llcell{b}{3.2,0}{2}
          \llcell[fill=green!25]{n}{1.6,-0.95}{5}
          \node[llvar] (cur) at (0,0.7) {cur};
          \node[llvar] (new) at (0.5,-0.95) {new};
          \draw[llptr] (cur) -- (av.north);
          \draw[llptr] (new) -- (nv.west);
          \draw[llptr] (an.center) -- (bv.west);
          \draw[llptr] (bn.center) -- ++(0.8,0) node[right, llvar] {...};
          \draw[llptr, red] (nn.center) -| (bv.south);
        \end{scope}
        % Step 2: the previous cell points at the new cell
        \begin{scope}[yshift=-3.1cm]
          \node[font=\scriptsize\bfseries, anchor=west] at (-0.5,1.25) {2.\ \texttt{cur\%next => new}};
          \llcell{a}{0,0}{1}
          \llcell{b}{3.2,0}{2}
          \llcell[fill=green!25]{n}{1.6,-0.95}{5}
          \node[llvar] (cur) at (0,0.7) {cur};
          \node[llvar] (new) at (0.5,-0.95) {new};
          \draw[llptr] (cur) -- (av.north);
          \draw[llptr] (new) -- (nv.west);
          \draw[llptr, gray!60, dashed] (an.center) -- (bv.west);
          \draw[llptr] (bn.center) -- ++(0.8,0) node[right, llvar] {...};
          \draw[llptr] (nn.center) -| (bv.south);
          \draw[llptr, red] (an.center) -| (nv.north);
        \end{scope}
      \end{tikzpicture}
    \end{column}

    \begin{column}{0.52\textwidth}
      \begin{itemize}
        \item To insert a cell after the one \texttt{cur} points to, allocate it and change two pointers.
        \item The order matters. Step 1 copies the link to the rest of the list into the new cell before step 2 overwrites it. In the other order, that link is lost: a memory leak.
        \item No other cell is touched or moved, however long the list is.
      \end{itemize}
      \lstinputlisting[firstline=25,lastline=29]{linked_list.f90}
    \end{column}
  \end{columns}
\end{frame}


\begin{frame}[fragile]{Linked lists: walking through a list}
  \centering
  \begin{tikzpicture}
    \node[llnext, label={[llvar]above:head}] (h) at (0,0) {};
    \fill (h.center) circle (1.3pt);
    \llcell{a}{1.7,0}{1}
    \llcell{b}{3.8,0}{5}
    \llcell{c}{5.9,0}{2}
    \llcell{d}{8.0,0}{3}
    \node[llvar, anchor=west] (nil) at (9.7,0) {null()};
    \draw[llptr] (h.center) -- (av.west);
    \draw[llptr] (an.center) -- (bv.west);
    \draw[llptr] (bn.center) -- (cv.west);
    \draw[llptr] (cn.center) -- (dv.west);
    \draw[llptr] (dn.center) -- (nil.west);
    % cur visits the cells one after another
    \node[llvar, gray] (cura) at ($(av.south)+(0,-0.8)$) {cur};
    \node[llvar, text=red] (curb) at ($(bv.south)+(0,-0.8)$) {cur};
    \node[llvar, gray] (curc) at ($(cv.south)+(0,-0.8)$) {cur};
    \node[llvar, gray] (curd) at ($(dv.south)+(0,-0.8)$) {cur};
    \node[llvar, gray] (cure) at (curd -| nil.south) {cur};
    \draw[llptr, gray, dashed] (cura) -- (av.south);
    \draw[llptr, red] (curb) -- (bv.south);
    \draw[llptr, gray, dashed] (curc) -- (cv.south);
    \draw[llptr, gray, dashed] (curd) -- (dv.south);
    \draw[llptr, gray, dashed] (cure) -- (nil.south);
    \draw[llptr, gray] (cura) -- (curb);
    \draw[llptr, gray] (curb) -- node[below=3pt, llvar, text=black] {cur => cur\%next} (curc);
    \draw[llptr, gray] (curc) -- (curd);
    \draw[llptr, gray] (curd) -- (cure);
  \end{tikzpicture}

  \vspace{0.8em}
  \begin{columns}[T]
    \begin{column}{0.5\textwidth}
      \begin{itemize}
        \item Start with \texttt{cur} pointing at the first cell.
        \item \texttt{cur => cur\%next} moves \texttt{cur} one cell forward.
        \item After the last cell, \texttt{cur} is null, \texttt{associated(cur)} is false, and the loop ends.
      \end{itemize}
    \end{column}

    \begin{column}{0.5\textwidth}
      \lstinputlisting[firstline=32,lastline=36]{linked_list.f90}
      This prints 1, 5, 2, and 3.
    \end{column}
  \end{columns}
\end{frame}


\begin{frame}[fragile]{Linked lists: removing a cell}
  \begin{columns}[T]
    \begin{column}{0.48\textwidth}
      \centering
      \begin{tikzpicture}
        % Step 1: the list bypasses the cell
        \begin{scope}
          \node[font=\scriptsize\bfseries, anchor=west] at (-0.5,1.25) {1.\ \texttt{cur\%next => old\%next}};
          \llcell{a}{0,0}{1}
          \llcell{b}{3.2,0}{2}
          \llcell[fill=red!20]{o}{1.6,-0.95}{5}
          \node[llvar] (cur) at (0,0.7) {cur};
          \node[llvar] (old) at (0.5,-0.95) {old};
          \draw[llptr] (cur) -- (av.north);
          \draw[llptr] (old) -- (ov.west);
          \draw[llptr, gray!60, dashed] (an.center) -| (ov.north);
          \draw[llptr] (on.center) -| (bv.south);
          \draw[llptr] (bn.center) -- ++(0.8,0) node[right, llvar] {...};
          \draw[llptr, red] (an.center) -- (bv.west);
        \end{scope}
        % Step 2: the cell is freed
        \begin{scope}[yshift=-3.1cm]
          \node[font=\scriptsize\bfseries, anchor=west] at (-0.5,1.25) {2.\ \texttt{deallocate(old)}};
          \llcell{a}{0,0}{1}
          \llcell{b}{3.2,0}{2}
          \llcell[fill=white, dashed, draw=gray!60, text=gray!60]{o}{1.6,-0.95}{5}
          \draw[red, thick] (ov.south west) -- (on.north east) (ov.north west) -- (on.south east);
          \node[llvar] (cur) at (0,0.7) {cur};
          \draw[llptr] (cur) -- (av.north);
          \draw[llptr] (an.center) -- (bv.west);
          \draw[llptr] (bn.center) -- ++(0.8,0) node[right, llvar] {...};
        \end{scope}
      \end{tikzpicture}
    \end{column}

    \begin{column}{0.52\textwidth}
      \begin{itemize}
        \item Keep a pointer to the cell (\texttt{old}), let the list bypass it, and then free it.
        \item Without \texttt{old}, nothing would point to the cell after step 1: a memory leak.
        \item To free a whole list, remove the first cell until \texttt{head} is null.
      \end{itemize}
      \lstinputlisting[firstline=40,lastline=42,belowskip=0pt]{linked_list.f90}
      \lstinputlisting[firstline=45,lastline=49,aboveskip=4pt]{linked_list.f90}
    \end{column}
  \end{columns}
\end{frame}


% \begin{frame}[fragile]{Pointers and recursive data structures}
%   \begin{columns}[T]
%     \column{0.52\textwidth}
%     \begin{block}{Linked lists via pointers}
%       \begin{itemize}
%         \item Derived types with pointer components support lists/trees.
%         \item Note subtle assignment semantics with pointer components:
%       \end{itemize}
% \begin{lstlisting}[language=Fortran]
% type(cell)        :: a
% type(cell),target :: b
% a = b
% ! equivalent to:
% ! a%val  = b%val
% ! a%next => b%next
% \end{lstlisting}
%       The pointer component is pointer-assigned (\texttt{=>}), not value-copied.
%     \end{block}
%
%     \column{0.48\textwidth}
% \begin{lstlisting}[language=Fortran]
% type cell
%   integer :: val
%   type(cell), pointer :: next
% end type cell
% \end{lstlisting}
%   \end{columns}
% \end{frame}
%
% \begin{frame}[fragile]{Linked list: construction}
%   \begin{columns}[T]
%     \column{0.52\textwidth}
%     \begin{block}{Idea}
%       Build a list with a \texttt{head} node, appending values read from input until termination. Pointers avoid copying and give O(1) append by keeping a tail pointer.
%     \end{block}
%
%     \column{0.48\textwidth}
% \begin{lstlisting}[language=Fortran]
% program thingy
%   implicit none
%   type(cell), target  :: head
%   type(cell), pointer :: curr, temp
%   integer :: k
%
%   head%val = 0
%   nullify(head%next)
%   curr => head
%   do
%     read*, k
%     allocate(temp)
%     temp%val = k
%     nullify(temp%next)
%     curr%next => temp
%     curr => temp
%   end do
% end program thingy
% \end{lstlisting}
%   \end{columns}
% \end{frame}
%
% \begin{frame}[fragile]{Linked list: traversal}
%   \begin{columns}[T]
%     \column{0.52\textwidth}
%     \begin{block}{Walk-through}
%       Traverse from \texttt{head} to the end by following \texttt{next}. Stop when \texttt{next} is disassociated.
%     \end{block}
%
%     \column{0.48\textwidth}
% \begin{lstlisting}[language=Fortran]
% curr => head
% do
%   print *, curr%val
%   if (.not. associated(curr%next)) exit
%   curr => curr%next
% end do
% \end{lstlisting}
%   \end{columns}
% \end{frame}



\begin{frame}[fragile]{Arrays of pointers}
  \begin{columns}[T]
    \column{0.52\textwidth}
    \begin{block}{Notes}
      \begin{itemize}
        \item Fortran has no arrays of pointers: \texttt{real, pointer :: p(:)} is one pointer to an array.
        \item To get an array of pointers, wrap a pointer in a derived type and make an array of that type.
        \item You cannot reference the pointer component of the whole array in one go.
        \item If desired, the container array could be \texttt{allocatable}.
      \end{itemize}
    \end{block}

    \column{0.48\textwidth}
\begin{lstlisting}[language=Fortran]
type iptr
  integer, pointer :: compon => null()
end type iptr

type(iptr), dimension(100) :: ints
integer, target :: k = 5

ints(10)%compon => k        ! ok
print *, ints(10)%compon    ! ok: 5
! print *, ints(:)%compon   ! not allowed
\end{lstlisting}
  \end{columns}
\end{frame}


\section{Custom Operators}

\begin{frame}
	\centering
	\Huge \textbf{Custom Operators}
\end{frame}

\begin{frame}[fragile]
  \frametitle{Arithmetic with derived types}
  \begin{columns}[T]
    \begin{column}{0.5\textwidth}
      \begin{itemize}
        \item Fortran allows you to define complex derived types.
        \item But using them directly for calculations leads to cumbersome code (and bugs, lots and lots of bugs).
        \item Here \texttt{c = s*a + b} has to be written out component by component.
      \end{itemize}
    \end{column}

    \begin{column}{0.5\textwidth}
    \lstinputlisting[firstline=1,lastline=6]{vec_components.f90}
    \lstinputlisting[firstline=8,lastline=17]{vec_components.f90}
    \end{column}
  \end{columns}
\end{frame}

\begin{frame}[fragile]
  \frametitle{Step 1: functions}
   Code can be simplified (and made more robust) by using functions.
  \begin{columns}[T]
    \begin{column}{0.5\textwidth}
          \lstinputlisting[firstline=1,lastline=18]{vec_functions.f90}
    \end{column}

    \begin{column}{0.5\textwidth}
    \lstinputlisting[firstline=20,lastline=29]{vec_functions.f90}
    \end{column}
  \end{columns}
\end{frame}

\begin{frame}[fragile]
  \frametitle{Step 2: operators}
  \begin{columns}[T]
    \begin{column}{0.5\textwidth}
          \lstinputlisting[firstline=1,lastline=22,aboveskip=0pt,belowskip=0pt]{vec_operators.f90}
    \end{column}

    \begin{column}{0.5\textwidth}
    \lstinputlisting[firstline=24,lastline=33,aboveskip=0pt]{vec_operators.f90}
    And defining custom operators can make your code even cleaner: the program now says \mbox{\texttt{c = s*a + b}}.
    \end{column}
  \end{columns}
\end{frame}


\begin{frame}[fragile]
  \frametitle{How to Define a Custom Operator}
  Custom operators are defined inside a module with an \texttt{interface operator} block.
  \vspace*{2mm}

  \begin{columns}[T]
    \begin{column}{0.5\textwidth}
      \begin{itemize}
            \item Standard operators (+, -, *, /)
      \end{itemize}
\begin{lstlisting}
interface operator(+)
   module procedure newsum
end interface
\end{lstlisting}
    \end{column}

    \begin{column}{0.5\textwidth}
      \begin{itemize}
            \item A new operator named \texttt{.op.}
      \end{itemize}
\begin{lstlisting}
interface operator(.op.)
   module procedure newop
end interface
\end{lstlisting}
    \end{column}
  \end{columns}

  \vspace*{2mm}
  The functions implementing the operators are defined later in the module after the \texttt{contains} keyword.
\end{frame}


\begin{frame}[fragile]
  \frametitle{Rules for custom operators}
  \begin{columns}[T]
    \begin{column}{0.48\textwidth}
      \begin{itemize}
        \item An operator function takes one argument (unary) or two (binary), all \texttt{intent(in)}.
        \item You can extend \texttt{+}, \texttt{-}, \texttt{*}, and \texttt{/} to your own types, but you cannot change what they do for the intrinsic types.
        \item Extended standard operators keep their usual precedence, so \texttt{s*a + b} works as expected.
        \item A named operator such as \texttt{.dot.} has the lowest precedence of all. Use parentheses.
      \end{itemize}
    \end{column}

    \begin{column}{0.52\textwidth}
      \lstinputlisting[firstline=64,lastline=74]{vec_dot.f90}
      The full module, with \texttt{+}, \texttt{-}, \texttt{*}, \texttt{/}, and \texttt{.dot.}, is in \texttt{examples/ex07\_operators}.
    \end{column}
  \end{columns}
\end{frame}


\section{Exercises}

\begin{frame}{Exercises: a star catalogue}
  \begin{block}{Tasks}
    \begin{enumerate}
      \item Write a module with a type \texttt{vector3d} and a type \texttt{star} that holds a name, a mass, and a position (a \texttt{vector3d}).
      \item In a main program, create an array of three stars with constructors, and print the name of each star.
      \item Print the total mass of the stars. Hint: \texttt{stars\%mass} is an array.
      \item Add a function \texttt{centre\_of\_mass(stars)} to the module that returns a \texttt{vector3d}, and test it.
      \item Try to print one whole \texttt{star} with \texttt{print *}. Then make the name \texttt{character(len=:), allocatable} and try again. What changes?
    \end{enumerate}
  \end{block}
\end{frame}


\begin{frame}{Exercises: vector operators}
  \begin{block}{Tasks}
    \begin{enumerate}
      \item Write a module with a type \texttt{vector3d} and the operators \texttt{+} and \texttt{-} for two vectors, and \texttt{*} for a real times a vector.
      \item Add the named operators \texttt{.dot.} and \texttt{.cross.} for the scalar and vector products.
      \item Test them: compute \mbox{\texttt{c = a .cross. b}} for two vectors of your choice, and check that \mbox{\texttt{c .dot. a}} and \mbox{\texttt{c .dot. b}} are both zero.
      \item What does \mbox{\texttt{2.0 * a .dot. b}} compute? And \mbox{\texttt{1.0 + a .dot. b}}? Explain both.
      \item Make all the operator functions \texttt{pure}. Why is that a sensible default?
    \end{enumerate}
  \end{block}
\end{frame}


\end{document}
